\begin{Exercise}[label={ex:trova_k}]
    Scrivere una funzione che, ricevendo in ingresso una lista dinamica e un intero 
    $k > 0$, individui il nodo che si trova k posizioni dopo il nodo centrale 
    e ne restituisca il puntatore al chiamante. 
    Se tale nodo non esiste, ad esempio a causa di un’insufficiente numero di nodi 
    nella lista, la funzione deve restituire il valore NULL.
    
    \textbf{BONUS}: scrivere una soluzione che prevede una sola scansione 
    della lista (una scansione intera più una scansione anche solo parziale 
    è considerata equivalente a due scansioni).
\end{Exercise}

\begin{Answer}[ref={ex:trova_k}]
  	\lstinputlisting[style=cstylesmall]{src/trova_k.c}
\end{Answer}